\reportheading{6.\quad Conclusions and design recommendation}
The single open doorway concentrated exchange at R1. The measured ceiling-temperature peak shifted towards the closed side. Wall heat transfer improved the high-power OP-6 prediction but deepened the low-power OP-1 cold bias. Agreement in peak position alone therefore gives an incomplete picture of temperature accuracy.

For the propylene design, we recommend 1.12~m$^3$/s with two ceiling exhausts, 0.24~m curtains and end openings preset to 0.44~m for OP-1 or 0.48~m for OP-6. This was the lowest-power feasible choice among the three tested flows. The 0.96~m$^3$/s setting breached the height constraint under refinement. At 1.12~m$^3$/s, the smallest of the six refined-grid margins was 55.8~mm, at $x=6.1$~m and 113.0~s under OP-6 at soot yield 0.05. Estimated short-branch power fell from 10.32~W at 1.28~m$^3$/s to 6.91~W. Fan selection would also require full-system losses and a fan curve.

In the original layout, BO and GA found the least severe breach earlier in a replay of 16 designs. In the revised layout, a separate replay ranked the six scenarios within each flow. Checking OP-6 at baseline soot first already gave the fewest possible checks in this comparison. Using BO or GA to choose a design and then using the PINN to schedule its FDS runs remains a proposed workflow.
